%% ---------------------------------------------------------------------
\section{El método indirecto: mirar la jacobiana}
\label{sec:ly-indirecto}
%% ---------------------------------------------------------------------

Cerca del origen, $f(x)=Ax+g(x)$ con
\begin{equation}
  A=\jacat fx{x=0},\qquad \frac{\norm{g(x)}}{\norm x}\to0\ \text{cuando}\ \norm x\to0 ,
  \label{eq:ly-linealizacion}
\end{equation}
si $f$ es $\Cont1$. La aproximación lineal $\dot x=Ax$ es la que se usó en los capítulos~\ref{cap:unidimensionales}, \ref{cap:bidimensionales} y~\ref{cap:orbitas}. Una matriz es \clave{Hurwitz} si todos sus valores propios tienen parte real negativa.

\begin{teorema}[label=thm:ly-indirecto]{Método indirecto de Lyapunov \citep[Teo.~4.7]{Kha02}}
Sea $x=0$ un equilibrio de \eqref{eq:ly-sistema}, con $f$ de clase $\Cont1$ cerca del origen, y $A$ su jacobiana \eqref{eq:ly-linealizacion}.
\begin{enumerate}
  \item Si $\operatorname{Re}\lambda_i<0$ para todos los valores propios de $A$ ($A$ Hurwitz), el origen es asintóticamente estable.
  \item Si $\operatorname{Re}\lambda_i>0$ para al menos un valor propio de $A$, el origen es inestable.
\end{enumerate}
\end{teorema}

La prueba de la parte~1 es el primer uso serio del método directo, y se hace en la \cref{sec:ly-lineales} (\cref{cor:ly-indirecto}): la función $V=x\T Px$ que certifica la estabilidad de $\dot x=Ax$ sirve también para $\dot x=Ax+g(x)$ cerca del origen. La parte~2 se prueba con el teorema de inestabilidad de Chetaev, que se estudia con los sistemas variantes en el tiempo (capítulo~9); la \cref{sec:ly-lineales} muestra de dónde sale la función que usa.

\begin{corolario}[label=cor:ly-exponencial]{Estabilidad exponencial local \citep[Cor.~4.3]{Kha02}}
El origen de \eqref{eq:ly-sistema} es exponencialmente estable si y solo si $A$ es Hurwitz.
\end{corolario}

\begin{ejemplo}[label=ej:ly-jacpendulo]{Los equilibrios del péndulo}
Clasifique los equilibrios $(0,0)$ y $(\pi,0)$ de \eqref{eq:ly-pendulo} con $a=1$ y $b=0.5$.
\solucion La jacobiana en $(x_1,0)$ es $\begin{bmatrix}0&1\\-a\cos x_1&-b\end{bmatrix}$. En $(0,0)$, $\lambda^2+b\lambda+a=0$ da $\lambda=-0.25\pm0.968j$: Hurwitz, y el origen es exponencialmente estable. En $(\pi,0)$, $\lambda^2+b\lambda-a=0$ da $\lambda=0.781$ y $-1.281$: una silla, inestable. Sin fricción ($b=0$), en el origen $\lambda=\pm j$: ninguno de los dos casos del teorema aplica.
\end{ejemplo}

\codigorepo{cap07_lyapunov/C7_Ej2_Indirecto.m}{jacobiano}{Jacobiana del péndulo en sus dos equilibrios}

\subsection{Cuando la jacobiana se calla}

El \cref{thm:ly-indirecto} no dice nada si $A$ tiene valores propios en el eje imaginario y ninguno en el semiplano derecho. No es una falla de la prueba: en ese caso los términos de orden superior deciden, y pueden decidir cualquier cosa.

\begin{ejemplo}[label=ej:ly-cubico]{Un caso que la linealización no decide}
Estudie el origen de $\dot x=ax-x^3$ en función de $a\in\R$.
\solucion La linealización es $\dot x=ax$. Si $a<0$, el origen es exponencialmente estable; si $a>0$, inestable (y aparecen dos equilibrios nuevos en $\pm\sqrt a$, con jacobiana $-2a<0$: la bifurcación de horquilla del \cref{cap:bifurcaciones}). Si $a=0$, la linealización es $\dot x=0$ y no decide. Las soluciones de $\dot x=-x^3$ son $x(t)=x_0/\sqrt{1+2x_0^2t}$: el origen es asintóticamente estable, pero no exponencialmente. Si el término cúbico tuviera el signo opuesto, $\dot x=+x^3$, la linealización sería la misma y el origen sería inestable (con escape en tiempo finito, como en la \cref{sec:un-escape}).
\end{ejemplo}

\begin{figure}[ht]
\centering
\includegraphics[width=\linewidth]{cap07/f_indirecto}
\caption{Izquierda: $\dot x=ax-x^3$ para $a=-1,0,1$; las tres tienen la misma forma lejos del origen y la linealización solo ve la pendiente en $0$. Derecha: con $a=0$ la convergencia es algebraica, no exponencial.}
\label{fig:ly-indirecto}
\end{figure}

\codigorepo{cap07_lyapunov/C7_Ej2_Indirecto.m}{cubico}{$\dot x=ax-x^3$ y la convergencia lenta de $\dot x=-x^3$}

\begin{trampa}[label=tr:ly-imaginario]{``Valores propios en el eje imaginario: el equilibrio es estable (pero no asintóticamente)''}
\plausible{Para un sistema lineal, un par $\pm j\omega$ simple da oscilaciones de amplitud constante; un sistema no lineal cerca del equilibrio se parece a su linealización, así que también oscila sin crecer.}
El parecido entre un sistema y su linealización solo está garantizado cuando la linealización es \emph{hiperbólica} (ningún valor propio con parte real nula). En el eje imaginario manda lo que la linealización desprecia. El péndulo sin fricción y $\dot x_1=x_2$, $\dot x_2=-x_1-x_2^3$ tienen la misma linealización ($\lambda=\pm j$): el primero es estable sin ser asintóticamente estable; el segundo es asintóticamente estable (con $V=\tfrac12(x_1^2+x_2^2)$, $\Vdot=-x_2^4\le0$, y el principio de invarianza del capítulo~8 completa la prueba). Y $\dot x_1=x_2$, $\dot x_2=-x_1+x_2^3$ es inestable.
\correcto{Con valores propios en el eje imaginario, el método indirecto no concluye nada. Hay que usar el método directo (o la teoría de la variedad central).}
\end{trampa}

\begin{pausa}
El \cref{thm:ly-indirecto} da estabilidad \emph{local}. ¿Qué tan grande puede ser la región donde vale la conclusión? ¿Lo dice el teorema?
\respuesta{No lo dice: solo garantiza que existe \emph{alguna} vecindad. Para el péndulo con fricción la región de atracción es enorme, pero no es todo el plano (las soluciones que dan la vuelta terminan en $(2\pi,0)$). Estimar esa región exige una función de Lyapunov explícita (\cref{sec:ly-region}).}
\end{pausa}
